\documentclass[11pt,professionalfont,aspectratio=43]{beamer}

% Assets
\usepackage[czech]{babel}				% Jazyk
\usepackage[a-2u]{pdfx}				    % Kopírování z pdfka
\usepackage{tikz}						% Schémata automatů
\usetikzlibrary{shadows, shadows.blur, shapes.geometric, arrows, positioning, overlay-beamer-styles, patterns, shadings, shapes.callouts}
% \usetikzlibrary{arrows}
\usepackage{pgfplots}
% \pgfplotsset{compat=1.15}
% \usepackage{mathrsfs}
\usepackage[export]{adjustbox}

\usepackage[T1]{fontenc}
\usefonttheme[onlymath]{serif}
% \usepackage{newtxmath}
\usepackage{pxfonts}
\usepackage{slantsc}
% \usepackage{varwidth}

\usepackage[linesnumbered,ruled,vlined]{algorithm2e}                % Pseudokód algoritmů
\usepackage[utf8]{inputenc}
% \usepackage{textcomp}
\usepackage{hyperref}
\usepackage{fancybox}

%%% FONTS
% \usefonttheme{serif}
% \usepackage{lmodern}
\usepackage{helvet}

\hypersetup{%
    pdfencoding=auto,
    pdfauthor={\insertauthor},
    pdftitle={\insertsubtitle}
}
\usepackage{float}
\usepackage{csquotes}					% české uvozovky
\usepackage{enumerate}					% enumerate environment
\usepackage{indentfirst}
\usepackage{tcolorbox}                  % Fancy boxíky
\usepackage{mathtools}
\usepackage{pifont}
\usepackage{cancel}
\usepackage{subcaption}
\usepackage{soul}
\usepackage{makecell}                   % Zalomení v buňce tabulky
\usepackage{xcolor}
\usepackage{graphicx}
\usepackage{tabularx}
\usepackage{amsmath}
\usepackage{amsthm}
\usepackage{colortbl}                   % Barvení buněk v tabulce
\usepackage{multicol}                   % Rozdělení obsahu do sloucpů

\usepackage{listings}                   % Úryvky z kódu
\definecolor{maroon}{rgb}{0.5,0,0}
\definecolor{forestgreen}{rgb}{0.13, 0.55, 0.13}

% \lstdefinelanguage{XML}
% {
%   basicstyle=\ttfamily,
%   morestring=[s]{"}{"},
%   morecomment=[s]{?}{?},
%   morecomment=[s]{!--}{--},
%   commentstyle=\color{forestgreen},
%   moredelim=[s][\color{black}]{>}{<},
%   moredelim=[s][\color{red}]{\ }{=},
%   stringstyle=\color{blue},
%   identifierstyle=\color{maroon}
% }
% Colors

\def\maincolR{0} % červená složka (0–255)
\def\maincolG{83} % zelená složka
\def\maincolB{120} % modrá složka
\definecolor{maincolor}{HTML}{8B00B5}

% \pgfmathtruncatemacro{\authorfooterbgR}{1*\maincolR}
% \pgfmathtruncatemacro{\authorfooterbgG}{1*\maincolG}
% \pgfmathtruncatemacro{\authorfooterbgB}{1*\maincolB}
% \definecolor{authorfooterbg}{RGB}{\authorfooterbgR,\authorfooterbgG,\authorfooterbgB}

% \pgfmathtruncatemacro{\titlefooterbgR}{0.722*\maincolR}   
% \pgfmathtruncatemacro{\titlefooterbgG}{0.974*\maincolG}   
% \pgfmathtruncatemacro{\titlefooterbgB}{0.739*\maincolB}
% \definecolor{titlefooterbg}{RGB}{\titlefooterbgR,\titlefooterbgG,\titlefooterbgB}

% Beamer theme
\colorlet{themecolor}{maincolor}
\colorlet{authorfooterbg}{maincolor}
\definecolor{titlefooterbg}{HTML}{9471B0}
\definecolor{titlefooterfg}{HTML}{000000}
\colorlet{titlecolorfg}{maincolor}
\definecolor{titleboxbg}{HTML}{E7E7E7}
\definecolor{datefooterbg}{HTML}{E3E3E3}
\colorlet{titleframeboxcolor}{maincolor}

% Background gradient
\definecolor{gradMiddle}{RGB}{212,248,255}
\definecolor{gradEnd}{RGB}{114,231,255}

% Block theme
\definecolor{blocktitlebg}{HTML}{F0DFA9}
\definecolor{blocktitlefg}{HTML}{957300}
\definecolor{blockbodybg}{HTML}{F5F0DF}

\definecolor{alertblockbodybg}{HTML}{FFE4E4}
\definecolor{alertblocktitlebg}{HTML}{ECA4A4}
\definecolor{alertblocktitlefg}{HTML}{B11313}

\definecolor{exampleblockbodybg}{HTML}{DAF0DA}
\definecolor{exampleblocktitlebg}{HTML}{BBE0BB}
\definecolor{exampleblocktitlefg}{HTML}{257A25}

\definecolor{mathfg}{HTML}{00185E}

\colorlet{bulletbg}{maincolor}
% \definecolor{bulletbg}{HTML}{}

% Beamer theme
\usetheme{Madrid}
\usecolortheme{seahorse}
% \useoutertheme[subsection=false]{miniframes}

\usecolortheme[named=themecolor]{structure}
\setbeamertemplate{frame numbering}[fraction]
\setbeamertemplate{navigation symbols}{}
\setbeamercovered{invisible}

% Header/footer
\setbeamercolor{author in head/foot}{bg=authorfooterbg,fg=white}
\setbeamercolor{title in head/foot}{bg=titlefooterbg,fg=titlefooterfg}
\setbeamercolor{date in head/foot}{bg=datefooterbg,fg=black}
% \setbeamercolor{title}{bg=datefooterbg}
% \setbeamercolor{title}{fg=titlecolorfg}
\setbeamerfont{title}{series=\bfseries}
\setbeamercolor{frametitle}{bg=datefooterbg}
% \setbeamercolor{section in head/foot}{bg=red,fg=cyan}
% \setbeamercolor{subsection in head/foot}{bg=blue,fg=yellow}

% Table of contents
\setbeamercolor{section in toc}{fg=black}
\setbeamercolor{subsection in toc}{fg=red}
\setbeamercolor{section number projected}{fg=black}

% Blocks
\setbeamertemplate{blocks}[rounded][shadow=true]

\setbeamercolor{block title}{bg=blocktitlebg, fg=blocktitlefg}
\setbeamercolor{block body}{parent=normal text,use=block title,bg=blockbodybg}

\setbeamercolor{block body alerted}{bg=alertblockbodybg}
\setbeamercolor{block title alerted}{bg=alertblocktitlebg,fg=alertblocktitlefg}

\setbeamercolor{block body example}{bg=exampleblockbodybg}
\setbeamercolor{block title example}{bg=exampleblocktitlebg,fg=exampleblocktitlefg}

% Math
\setbeamercolor{math text}{fg=mathfg}

% Itemize a enumerate
\setbeamertemplate{itemize items}[circle]
\setbeamertemplate{enumerate items}[circle]
\setbeamercolor{itemize item}{fg=bulletbg,bg=white}
\setbeamercolor{itemize subitem}{fg=bulletbg,bg=white}
\setbeamercolor{itemize subsubitem}{fg=bulletbg,bg=white}

\setbeamertemplate{enumerate items}[circle]
\setbeamercolor{item projected}{bg=bulletbg}

% Titlepage
\defbeamertemplate*{title page}{customized}[1][]
{
    \vbox{}%
    \vfill%
    \begin{centering}
        \begin{tikzpicture}
            \node[thick, rectangle, rounded corners=5mm, drop shadow={shadow xshift=0mm, shadow yshift=2mm, opacity=0.2, blur shadow, shadow blur steps=10}, inner sep=5mm, fill=titleboxbg, fill opacity=1] {%
                {\usebeamerfont{title}\textcolor{titlecolorfg}{\inserttitle}\par}%
            };
        \end{tikzpicture}
        \vskip1em\par%
        \ifx\insertsubtitle\@empty%
        \else%
            {\usebeamerfont{subtitle}\insertsubtitle\par}%
        \fi%
        \vskip1em\par%
        \ifx\insertauthor\@empty%
        \else%
            {\usebeamerfont{author}\insertauthor}\\%
            \texttt{\email}
        \fi
        \vskip0.3em\par%
        % \ifx\insertinstitute\@empty%
        % \else%
        %     {\usebeamerfont{institute}\insertinstitute\par}%
        % \fi%
        \vskip1em\par%
        \ifx\inserttitlegraphic\@empty%
        \else%
            \vskip1em\par%
            \inserttitlegraphic\par%
        \fi%
        \vskip1em\par%
        \ifx\insertdate\@empty%
        \else%
            {\usebeamerfont{date}\insertdate\par}%
        \fi%
    \end{centering}%
    \vfill%
}
% Enumerate
%\setlist[enumerate]{topsep=0pt,itemsep=-1ex,partopsep=1ex,parsep=1ex,label=(\arabic*)}

\MakeOuterQuote{"}

%%% Makra pro definice, věty, tvrzení, příklady, ... (vyžaduje baliček amsthm)

\theoremstyle{plain}
\newtheorem{thm}{Věta}[section]
\newtheorem{lem}[theorem]{Lemma}
\newtheorem{prop}[theorem]{Tvrzení}
\newtheorem{cor}[theorem]{Důsledek}
\newtheorem*{prop*}{Tvrzení}

\theoremstyle{definition}
\newtheorem{define}[theorem]{Definice}
% \newtheorem{example}[theorem]{Příklad}
% \newtheorem{remark}[theorem]{Poznámka}
% \newtheorem{convention}[theorem]{Úmluva}
% \newtheorem{denoting}[theorem]{Značení}

\newenvironment{defblock}[1][]{
  \begin{block}{\ifthenelse{\equal{#1}{}}
  {Definice}
  {Definice (#1)}}
}{\end{block}}
\newenvironment{thmblock}[1][]{
  \begin{block}{\ifthenelse{\equal{#1}{}}
  {Věta}
  {Věta (#1)}}
}{\end{block}}
\newenvironment{corblock}[1][]{
  \begin{block}{\ifthenelse{\equal{#1}{}}
  {Důsledek}
  {Důsledek (#1)}}
}{\end{block}}
\newenvironment{propblock}[1][]{
  \begin{block}{\ifthenelse{\equal{#1}{}}
  {Tvrzení}
  {Tvrzení (#1)}}
}{\end{block}}
\newenvironment{exblock}[1][]{
  \begin{exampleblock}{\ifthenelse{\equal{#1}{}}
  {Příklad}
  {Příklad (#1)}}
}{\end{exampleblock}}


% Colors
\definecolor{lightblue}{HTML}{009AD4}
\definecolor{lightgreen}{HTML}{68FF00}
\definecolor{darkgreen}{HTML}{00B600}
\definecolor{darkblue}{HTML}{0265b3}
\definecolor{darkred}{HTML}{DA0707}
\definecolor{lightred}{HTML}{FF5100}
\definecolor{orange}{HTML}{FFE000}
\definecolor{codeblue}{HTML}{007eed}
\definecolor{codegreen}{rgb}{0,0.6,0}
\definecolor{codegray}{rgb}{0.5,0.5,0.5}
\definecolor{codebeige}{HTML}{D4A000}
\definecolor{codepink}{HTML}{FF0055}
\definecolor{backcolour}{rgb}{0.95,0.95,0.92}

\definecolor{decproblemtitlefg}{HTML}{957300}
\definecolor{decproblembg}{HTML}{F5F0DF}

\definecolor{timportantboxcolor}{HTML}{A97100}
\definecolor{talertboxcolor}{HTML}{B70000}
\definecolor{tnoticeboxcolor}{HTML}{159C00}
\definecolor{tinfoboxcolor}{HTML}{004DFF}

\newcommand{\markred}[1]{\textcolor{lightred}{#1}}
\newcommand{\markdarkred}[1]{\textcolor{darkred}{#1}}
\newcommand{\markgreen}[1]{\textcolor{lightgreen}{#1}}
\newcommand{\markdarkgreen}[1]{\textcolor{darkgreen}{#1}}
\newcommand{\markorange}[1]{\textcolor{orange}{#1}}
\newcommand{\markblue}[1]{\textcolor{lightblue}{#1}}
\newcommand{\markdarkblue}[1]{\textcolor{darkblue}{#1}}

% Math mode settings
% \AtBeginDocument{
%     \setlength{\abovedisplayskip}{3pt}
%     \setlength{\belowdisplayskip}{-4pt}
% }

% Inline images
\newcommand{\inlineimgscale}{1.1}

% X and check mark
\newcommand{\cmark}{\markgreen{\ding{51}}}
\newcommand{\xmark}{\markred{\ding{55}}}

% Math
\newcommand{\R}{\mathbb{R}}
\newcommand{\C}{\mathbb{C}}
\newcommand{\N}{\mathbb{N}}
\newcommand{\Z}{\mathbb{Z}}
\newcommand{\Q}{\mathbb{Q}}

\newcommand{\T}[1]{#1^\top}                     % Tranpozice matice
\newcommand{\compl}[1]{#1^\complement}
\newcommand{\mapping}[3]{#1:#2 \rightarrow #3}
\newcommand{\mapto}[3]{#1:#2 \mapsto #3}
\newcommand{\set}[1]{\left\{#1\right\}}
\newcommand{\powset}{\mathcal{P}}
\newcommand{\code}[1]{\langle#1\rangle}
\DeclareMathOperator*{\argmin}{arg\,min}
\DeclareMathOperator*{\argmax}{arg\,max}

% Statistika a AI
\newcommand{\average}[1]{\overline{#1}}
\DeclareMathOperator{\median}{med}
\DeclareMathOperator{\Var}{var}
\DeclareMathOperator{\Cov}{cov}
\DeclareMathOperator{\MSE}{MSE}
\DeclareMathOperator{\SSE}{SSE}

\DeclareMathOperator{\TP}{TP}
\DeclareMathOperator{\TN}{TN}
\DeclareMathOperator{\FP}{FP}
\DeclareMathOperator{\FN}{FN}
\DeclareMathOperator{\Accuracy}{Acc}
\DeclareMathOperator{\Precision}{Prec}
\DeclareMathOperator{\Recall}{Rec}
\DeclareMathOperator{\Fscore}{F}
\DeclareMathOperator{\Gini}{Gini}

\DeclareMathOperator{\mspan}{span}
\DeclareMathOperator{\mrank}{rank}
\DeclareMathOperator{\tg}{tg}
\DeclareMathOperator{\id}{id}
\DeclareMathOperator{\dom}{dom}
\DeclareMathOperator{\rng}{rng}
\renewcommand{\emptyset}{\varnothing}
\newcommand{\binary}[1]{\ensuremath{\left\langle #1\right\rangle}_B}

% DFS in and out lists
\DeclareMathOperator{\inlist}{in}
\DeclareMathOperator{\outlist}{out}
% \DeclareMathOperator{\lowlist}{low}
\DeclareMathOperator{\key}{key}

% A* macro
\newcommand{\Astar}{$\textcolor{black}{\textup{A}^{*}}$}

% Complexity classes
\DeclareMathOperator{\classP}{P}
\DeclareMathOperator{\classNP}{NP}
\DeclareMathOperator{\classTime}{TIME}
\DeclareMathOperator{\classNTime}{NTIME}
\DeclareMathOperator{\classSpace}{SPACE}
\DeclareMathOperator{\classNSpace}{NSPACE}

\DeclareMathOperator{\SAT}{SAT}
\DeclareMathOperator{\THREESAT}{3-SAT}
\DeclareMathOperator{\UNSAT}{UNSAT}
\DeclareMathOperator{\NZMNA}{NzMna}
\DeclareMathOperator{\KLIKA}{Klika}
\DeclareMathOperator{\ROZVRH}{Rozvrh}
\DeclareMathOperator{\SUDOKU}{Sudoku}
\newcommand{\HALT}{\ensuremath{\textsc{Halt}}}
\newcommand{\FACTOR}{\ensuremath{\textsc{Factor}}}

\DeclareMathOperator{\regex}{RegE}

%%% Barevné boxíky (notice / info / important / alert)
%
% Šířka boxu se nezadává ručně -- box se přizpůsobí šířce svého obsahu.
% Obsah se vysází do prostředí varwidth, takže krátký text zůstane na
% jednom řádku; delší text se zalomí, nejvýše však na šířku \linewidth
% (tj. \textwidth mimo sloupce a minipage). Pokud je titulek širší než
% obsah, řídí šířku boxu titulek.

\newsavebox{\tboxcontentbox}
\newsavebox{\tboxtitlebox}
\newlength{\tboxwidth}
\newlength{\tboxpadding}

% Geometrie boxu (musí odpovídat stylu tautobox níže):
% 2 * (boxrule + boxsep + left)
\setlength{\tboxpadding}{\dimexpr 0.5mm+1mm+4mm\relax}
\setlength{\tboxpadding}{2\tboxpadding}

\tcbset{tautobox/.style={boxrule=0.5mm, boxsep=1mm, left=4mm, right=4mm,
                         fonttitle=\bfseries}}

% Začátek sběru obsahu do boxu
\newcommand{\tboxcollect}{%
  \begin{lrbox}{\tboxcontentbox}%
    \begin{varwidth}{\dimexpr\linewidth-\tboxpadding\relax}%
}

% Ukončení sběru a výpočet výsledné šířky
% #1 = titulek (pro měření; prázdný u variant bez titulku)
\newcommand{\tboxmeasure}[1]{%
    \end{varwidth}%
  \end{lrbox}%
  \sbox{\tboxtitlebox}{\bfseries #1}%
  \setlength{\tboxwidth}{\wd\tboxcontentbox}%
  \ifdim\wd\tboxtitlebox>\tboxwidth
    \setlength{\tboxwidth}{\wd\tboxtitlebox}%
  \fi
  \addtolength{\tboxwidth}{\tboxpadding}%
  \ifdim\tboxwidth>\linewidth
    \setlength{\tboxwidth}{\linewidth}%
  \fi
}

% Vysázení boxu s titulkem: #1 = barva, #2 = titulek
\newcommand{\tboxrenderbox}[2]{%
  \tboxmeasure{#2}%
  \begin{center}%
    \begin{tcolorbox}[tautobox, colback=#1!15, colframe=#1,
                      title={#2}, width=\tboxwidth]
      \usebox{\tboxcontentbox}%
    \end{tcolorbox}%
  \end{center}%
}

% Vysázení boxu bez titulku: #1 = barva
\newcommand{\tboxrenderframe}[1]{%
  \tboxmeasure{}%
  \begin{center}%
    \begin{tcolorbox}[tautobox, colback=#1!15, colframe=#1,
                      width=\tboxwidth]
      \usebox{\tboxcontentbox}%
    \end{tcolorbox}%
  \end{center}%
}

% Poznámky (notice)
\newenvironment{tnoticebox}[1]{%
  \def\tboxtitletext{#1}\tboxcollect
}{%
  \tboxrenderbox{tnoticeboxcolor}{\tboxtitletext}%
}

\newenvironment{tnoticeframe}{%
  \tboxcollect
}{%
  \tboxrenderframe{tnoticeboxcolor}%
}

% Informační bloky (info)
\newenvironment{tinfobox}[1]{%
  \def\tboxtitletext{#1}\tboxcollect
}{%
  \tboxrenderbox{tinfoboxcolor}{\tboxtitletext}%
}

\newenvironment{tinfoframe}{%
  \tboxcollect
}{%
  \tboxrenderframe{tinfoboxcolor}%
}

% Důležité bloky (important)
\newenvironment{timportantbox}[1]{%
  \def\tboxtitletext{#1}\tboxcollect
}{%
  \tboxrenderbox{timportantboxcolor}{\tboxtitletext}%
}

\newenvironment{timportantframe}{%
  \tboxcollect
}{%
  \tboxrenderframe{timportantboxcolor}%
}

% Varovné bloky (alert)
\newenvironment{talertbox}[1]{%
  \def\tboxtitletext{#1}\tboxcollect
}{%
  \tboxrenderbox{talertboxcolor}{\tboxtitletext}%
}

\newenvironment{talertframe}{%
  \tboxcollect
}{%
  \tboxrenderframe{talertboxcolor}%
}

% Decision problem
\newcommand{\decproblem}[3]{
    \begin{center}
        \begin{minipage}{.9\textwidth}
            \begin{table}[!htbp]
                \centering
                \bgroup
                \def\arraystretch{1.8}
                \begin{tabularx}{\textwidth}{|rX|}
                \hline
                \rowcolor{decproblembg} \multicolumn{2}{|c|}{\textcolor{decproblemtitlefg}{\textsc{#1}}} \\ \hline
                \markdarkred{\textbf{Vstup:}} & #2                \\ \hline
                \rowcolor{decproblembg} \markdarkred{\textbf{Otázka:}} & #3               \\ \hline
                \end{tabularx}
                \egroup
            \end{table}
        \end{minipage}
    \end{center}
}
\newcommand{\bigO}{\mathcal{O}}

% TODO
\newcommand{\todo}[1]{\textcolor{red}{(\noindent TODO: #1.)}}


\newcommand{\hlcode}[1]{\colorbox{red}{#1}}

% Algorithms
\numberwithin{algocf}{section}
\renewcommand{\algorithmcfname}{Algoritmus}
\SetKwInput{KwIn}{Vstup}
\SetKwInput{KwOut}{Výstup}
\SetKwProg{Function}{Funkce}{}{end}
\setlength{\algomargin}{25pt}
\renewcommand{\thealgocf}{}     % Vypnutí číslování algoritmů

% Definice stylů pro uzly
\tikzstyle{startstop} = [rectangle, rounded corners, minimum width=1.5cm, minimum height=0.5cm,text centered, draw=black, fill=red!30]
\tikzstyle{process}   = [rectangle, minimum width=1.5cm, minimum height=0.5cm, text centered, draw=black, fill=orange!30]
\tikzstyle{decision}  = [diamond, aspect=2, text centered, draw=black, fill=green!30]
\tikzstyle{arrow}     = [thick,->,>=stealth]

\tikzstyle{vertex}    = [thick, draw, circle, minimum size=4mm, inner sep=1pt, outer sep=1pt, align=center]
\tikzstyle{vertexplain}    = [minimum size=4mm, inner sep=1pt, outer sep=1pt, align=center]
\tikzstyle{verteximg}    = [draw, circle, minimum size=4mm, inner sep=0pt, outer sep=1pt, align=center, clip]
\tikzstyle{verteximgplain}    = [minimum size=4mm, inner sep=0pt, outer sep=1pt, align=center]
\tikzstyle{task}    = [draw, rounded corners=2pt, minimum width=14mm, minimum height=8mm, inner sep=2pt, align=center]

\tikzstyle{edge}    = [thick,>=stealth]
\tikzstyle{diredge}    = [thick, ->, >=stealth]

%% Automaty
\tikzstyle{state}=[%
    draw,
    circle,
    thick,
    minimum size=10mm,
    inner sep=0pt,
    outer sep=0pt
]

\tikzstyle{acceptstate}=[%
    state,
    double,
    double distance=1pt
]

\tikzstyle{transition}=[%
    ->,
    >=stealth,
    thick
]

\newcommand<>{\state}[4][]{%
    \node#5[state,#1] (#2) at #3 {#4};
}

\newcommand<>{\acceptingstate}[4][]{%
    \node#5[acceptstate,#1] (#2) at #3 {#4};
}

\newcommand<>{\transitionarrow}[5][]{%
    \path#6 (#2) edge[transition,#1] node[#3] {#4} (#5);
}

\newcommand<>{\initialstate}[6][]{%
    \node#7[state,#1] (#2) at #3 {#4};
    \draw#7[transition] #5 -- (#2.#6);
}

%%%%%

% Frames
\newcommand{\tableofcontentsframe}{
    \begin{frame}{Obsah}
        \tableofcontents
    \end{frame}
}

% \definecolor{boxcolor}{HTML}{00820a}
\newcommand{\titleframe}[1]{%
  \begin{frame}[plain]
    \begin{tikzpicture}[remember picture, overlay]
      \node at (current page.center)
        [draw=titleframeboxcolor, line width=1pt,
         inner xsep=10mm, inner ysep=4mm, rounded corners=2mm,
         text width=0.8\textwidth, align=center,
         execute at begin node=\setlength{\baselineskip}{2.2em}]
        {\Huge\markred{#1}\normalfont};
    \end{tikzpicture}
  \end{frame}
  \addtocounter{framenumber}{-1}%
}

\newenvironment{titleframeenv}[1]{
    \begin{frame}[plain]
        \begin{tcolorbox}[colback=white,boxrule=1pt]
            \begin{center}
                \Huge\markred{#1}\normalfont
            \end{center}
        \end{tcolorbox}
}{\addtocounter{framenumber}{-1}\end{frame}}

\newcommand{\icon}[1]{
  \includegraphics[height=\fontcharht\font`\B]{#1}
}

% Obrázek, který se vždy vejde na snímek. Adjustbox jej v případě potřeby
% zmenší tak, aby nepřesáhl šířku textu ani zadanou část výšky snímku; díky
% tomu sedí sazba ve všech používaných poměrech stran (4:3, 16:9 i 16:10).
% Volitelným argumentem je maximální výška (výchozí 0.68\textheight), povinným
% cesta k souboru s obrázkem (TikZ zdroj vkládaný přes \input).
\newcommand{\obrazek}[2][0.68\textheight]{%
    \begin{figure}
        \centering
        \begin{adjustbox}{max width=\linewidth, max totalheight=#1, center}
            \input{#2}
        \end{adjustbox}
    \end{figure}%
}

% Totéž pro rastrové či vektorové obrázky vkládané přes \includegraphics.
\newcommand{\obrazekgrf}[2][0.68\textheight]{%
    \begin{figure}
        \centering
        \includegraphics[width=\linewidth, max totalheight=#1, center]{#2}
    \end{figure}%
}

% Závěrečný snímek se seznamem zdrojů. Prvním (volitelným) argumentem je název
% sekce i snímku, druhým čárkou oddělený seznam klíčů z ../assets/literatura.bib.
% Citace se sázejí podle ČSN ISO 690, v pořadí, v jakém jsou klíče uvedeny.
\newcommand{\zdrojeframe}[2][Zdroje]{%
    \section{#1}
    \nocite{#2}%
    \begin{frame}[allowframebreaks]{#1}
        \printbibliography[heading=none]
    \end{frame}
}
\input{../assets/listing.tex}

\def\presentationtitle{Kompilátory}

\begin{document}

    \title[\presentationtitle]{\textsc{Teoretická informatika}}
\subtitle{\presentationtitle}
\author{David Weber}
\def\email{weber3@spsejecna.cz}
\institute{SPŠE Ječná}
\date{%
    Rok 
    \pgfmathsetmacro{\firstyear}{ifthenelse(\month<9, \year-1, \year)}%
    \pgfmathsetmacro{\secondyear}{ifthenelse(\month<9, \year, \year+1)}%
    \pgfmathtruncatemacro{\secondyearlasttwo}{mod(\secondyear,100)}%
    \pgfmathprintnumber[fixed,precision=0,zerofill,1000 sep={}]{\firstyear}/%
    \pgfmathprintnumber[fixed,precision=0,zerofill,1000 sep={}]{\secondyearlasttwo}%
}
\titlegraphic{\includegraphics[width=0.25\textwidth]{../images/SPSE-Jecna_Logo.pdf}}

\begin{frame}[plain]
    \titlepage
\end{frame}

    \section{Kompilátory}

    \titleframe{Kompilátory}

    \subsection{Úvod}

    \begin{frame}{Co je kompilátor?}

        \begin{itemize}
            \visible<2->{\item \textbf{Kompilátor} je program, který překládá jiný program.}
            \visible<3->{\item Vstupem je obvykle \textbf{zdrojový kód}.}
            \visible<4->{\item Výstupem je \textbf{cílový kód} -- např. assembler, mezikód nebo strojový kód.}
        \end{itemize}

        \vspace{0.25cm}

        \centering
        \begin{tikzpicture}[
            box/.style={
                draw,
                rounded corners,
                thick,
                align=center,
                minimum width=2.6cm,
                minimum height=0.8cm
            },
            arrow/.style={->, thick},
            node distance=0.75cm
        ]

            \visible<3->{\node[box] (src) {zdrojový\\kód};}
            \visible<2->{\node[box, right=of src] (comp) {kompilátor};}
            \visible<4->{\node[box, right=of comp] (out) {cílový\\kód};}

            \visible<3->{\draw[arrow] (src) -- (comp);}
            \visible<4->{\draw[arrow] (comp) -- (out);}

            \visible<5->{
                \node[below=0.25cm of src, align=center] {
                    \scriptsize
                    \shortstack{
                        \texttt{C}\\[-0.5mm]
                        \texttt{Java}\\[-0.5mm]
                        \texttt{C\#}
                    }
                };

                \node[below=0.25cm of out, align=center] {
                    \scriptsize
                    \shortstack{
                        assembler\\[-0.5mm]
                        mezikód\\[-0.5mm]
                        strojový kód
                    }
                };
            }

        \end{tikzpicture}

        \vspace{0.3cm}

        \visible<6->{
            \begin{tinfoframe}[0.88\linewidth]
                Program, který napíše člověk, se musí převést do podoby,
                se kterou dokáže pracovat počítač nebo běhové prostředí.
            \end{tinfoframe}
        }

    \end{frame}

    \subsection{Od zdrojového kódu ke spustitelnému souboru}

    \begin{frame}{Od zdrojového kódu ke spustitelnému souboru}

        \begin{itemize}
            \visible<2->{\item Překlad programu často neprobíhá v jednom kroku.}
            \visible<3->{\item Například u jazyka C může mít překladový řetězec několik částí.}
        \end{itemize}

        \vspace{0.15cm}

        \centering
        \begin{tikzpicture}[
            file/.style={
                draw,
                rounded corners,
                thick,
                align=center,
                font=\scriptsize,
                minimum width=1.25cm,
                minimum height=0.65cm
            },
            tool/.style={
                draw,
                rounded corners,
                thick,
                align=center,
                font=\scriptsize,
                minimum width=1.55cm,
                minimum height=0.65cm
            },
            arrow/.style={->, thick}
        ]

            \node<3->[file] (c) at (0,0) {\texttt{hello.c}};
            \node<3->[tool] (pre) at (1.8,0) {preprocesor};
            \node<3->[file] (i) at (3.6,0) {\texttt{hello.i}};

            \draw<3->[arrow] (c) -- (pre);
            \draw<3->[arrow] (pre) -- (i);

            \node<4->[tool] (comp) at (5.4,0) {kompilátor};
            \node<4->[file] (s) at (7.2,0) {\texttt{hello.s}};

            \draw<4->[arrow] (i) -- (comp);
            \draw<4->[arrow] (comp) -- (s);

            \node<5->[tool] (asm) at (5.4,-1.25) {assembler};
            \node<5->[file] (o) at (3.6,-1.25) {\texttt{hello.o}};

            \draw<5->[arrow] (s.south) -- ++(0,-0.35) -- (asm.east);
            \draw<5->[arrow] (asm) -- (o);

            \node<6->[tool] (link) at (1.8,-1.25) {linker};
            \node<6->[file] (exe) at (0,-1.25) {\texttt{hello}};

            \draw<6->[arrow] (o) -- (link);
            \draw<6->[arrow] (link) -- (exe);

        \end{tikzpicture}

        \vspace{0.15cm}

        \begin{columns}[T,onlytextwidth]
            \begin{column}{0.48\textwidth}
                \begin{itemize}
                    \visible<3->{\item \textbf{Preprocesor}: zpracuje např. \texttt{\#include}.}
                    \visible<4->{\item \textbf{Kompilátor}: vytvoří assembler.}
                \end{itemize}
            \end{column}

            \begin{column}{0.48\textwidth}
                \begin{itemize}
                    \visible<5->{\item \textbf{Assembler}: vytvoří objektový soubor.}
                    \visible<6->{\item \textbf{Linker}: spojí části programu a knihovny.}
                \end{itemize}
            \end{column}
        \end{columns}

    \end{frame}

    \subsection{Mezikód a strojový kód}

    \begin{frame}{Mezikód a strojový kód}
        \pause
        \begin{columns}[T,onlytextwidth]
            \begin{column}{0.49\textwidth}
                \textbf{Strojový kód}
                \begin{itemize}
                    \visible<3->{\item je určen pro konkrétní procesor,}
                    \visible<4->{\item procesor jej umí přímo vykonat,}
                    \visible<5->{\item je závislý na dané platformě.}
                \end{itemize}
            \end{column}

            \begin{column}{0.49\textwidth}
                \textbf{Mezikód}
                \begin{itemize}
                    \visible<6->{\item je mezilehlá reprezentace programu,}
                    \visible<7->{\item zpracuje jej virtuální stroj nebo JIT,}
                    \visible<8->{\item usnadňuje přenositelnost programu.}
                \end{itemize}
            \end{column}
        \end{columns}

        \vspace{0.1cm}

        \centering
        \begin{tikzpicture}[
            box/.style={
                draw,
                rounded corners,
                thick,
                align=center,
                font=\scriptsize,
                minimum width=1.35cm,
                minimum height=0.6cm
            },
            widebox/.style={
                draw,
                rounded corners,
                thick,
                align=center,
                font=\scriptsize,
                minimum width=1.65cm,
                minimum height=0.6cm
            },
            arrow/.style={->, thick}
        ]

            \visible<4->{
                \node[box] (c) at (0,0) {\texttt{C/C++}};
                \node[widebox] (comp1) at (1.9,0) {kompilátor};
                \node[widebox] (mc1) at (3.9,0) {strojový\\kód};
                \node[box] (cpu1) at (5.8,0) {procesor};

                \draw[arrow] (c) -- (comp1);
                \draw[arrow] (comp1) -- (mc1);
                \draw[arrow] (mc1) -- (cpu1);

                \node[font=\scriptsize, align=right] at (-1.8,0) {přímá\\kompilace};
            }

            \visible<7->{
                \node[box] (java) at (0,-1.25) {\texttt{Java}\\\texttt{C\#}};
                \node[widebox] (comp2) at (1.9,-1.25) {kompilátor};
                \node[box] (ir) at (3.9,-1.25) {mezikód};
                \node[box] (jit) at (5.8,-1.25) {VM/JIT};
                \node[widebox] (mc2) at (7.65,-1.25) {strojový\\kód};

                \draw[arrow] (java) -- (comp2);
                \draw[arrow] (comp2) -- (ir);
                \draw[arrow] (ir) -- (jit);
                \draw[arrow] (jit) -- (mc2);

                \node[font=\scriptsize, align=right] at (-1.8,-1.25) {překlad přes\\mezikód};
            }

        \end{tikzpicture}

        \vspace{0.15cm}

        \visible<9->{
            \begin{tinfoframe}
                Mezikód umožňuje přeložit program do společné podoby,
                která se na konkrétním počítači převede až později.
            \end{tinfoframe}
        }

    \end{frame}

    \subsection{Fáze překladače}

    \begin{frame}{Fáze překladače}

        \begin{itemize}
            \visible<2->{\item Překladač nezpracovává program najednou.}
            \visible<3->{\item Zdrojový kód postupně převádí do jednodušších reprezentací.}
            \visible<4->{\item Jednotlivé fáze se obvykle dělí na}
            \visible<5->{
                \begin{itemize}
                    \item \textbf{přední část} -- závisí hlavně na zdrojovém jazyce,
                    \item \textbf{koncovou část} -- závisí hlavně na cílovém počítači.
                \end{itemize}
            }
        \end{itemize}

        \vspace{0.15cm}

        \centering
        \begin{tikzpicture}[
            phase/.style={
                draw,
                rounded corners,
                thick,
                align=center,
                font=\scriptsize,
                minimum width=1.75cm,
                minimum height=0.65cm
            },
            lang/.style={
                draw,
                rounded corners,
                thick,
                align=center,
                font=\scriptsize,
                minimum width=1.45cm,
                minimum height=0.65cm
            },
            arrow/.style={->, thick}
        ]
            % vnější obdélník
            \visible<2->{
                \draw[thick, rounded corners] (-1.3,1.5) rectangle (7.8,-2.5);
                \node[draw, fill=white, font=\small\bfseries, inner xsep=6pt, inner ysep=2pt]
                    at (3.125,1.5) {Překladač};
            }

            \node<2->[lang,fill=red!40!white]  (src) at (0,0) {zdrojový\\kód};
            \node<3->[phase,fill=red!40!white] (lex) at (2.05,0) {lexikální\\analýza};
            \node<4->[phase,fill=red!40!white] (syn) at (4.10,0) {syntaktická\\analýza};
            \node<5->[phase,fill=red!40!white] (sem) at (6.15,0) {sémantická\\analýza};

            \node<6->[lang,fill=orange!40!white]  (ir)  at (6.15,-1.15) {mezikód};
            \node<7->[phase,fill=orange!40!white] (opt) at (4.10,-1.15) {optimalizace};
            \node<8->[phase,fill=orange!40!white] (gen) at (2.05,-1.15) {generování\\kódu};
            \node<9->[lang,fill=orange!40!white]  (tgt) at (0,-1.15) {cílový\\kód};

            \draw<3->[arrow] (src) -- (lex);
            \draw<4->[arrow] (lex) -- (syn);
            \draw<5->[arrow] (syn) -- (sem);
            \draw<6->[arrow] (sem) -- (ir);
            \draw<7->[arrow] (ir) -- (opt);
            \draw<8->[arrow] (opt) -- (gen);
            \draw<9->[arrow] (gen) -- (tgt);

            \node<5->[font=\scriptsize] at (3.1,0.85) {\textbf{přední část (front end)}};
            \node<8->[font=\scriptsize] at (3.1,-2.0) {\textbf{koncová část (back end)}};

        \end{tikzpicture}

    \end{frame}

    \section{Lexikální analýza}

    \titleframe{Lexikální analýza}

    \subsection{Hlavní činnost}

    \begin{frame}{Lexikální analýza}

        \begin{itemize}
            \visible<2->{\item \textbf{Lexikální analýza} je první analytická fáze překladače, kterou zařizuje tzv. \textbf{lexer}.}
            \visible<3->{\item Vstupem je posloupnost znaků zdrojového programu.}
            \visible<4->{\item Výstupem je posloupnost \textbf{lexikálních symbolů}.}
            \visible<5->{\item Těmto symbolům se také říká \textbf{tokeny}.}
        \end{itemize}

        \vspace{0.15cm}

        \centering
        \begin{tikzpicture}[
            box/.style={
                draw,
                rounded corners,
                thick,
                align=center,
                font=\scriptsize,
                minimum width=2.15cm,
                minimum height=0.75cm
            },
            arrow/.style={->, thick}
        ]

            \node<3->[box] (src) at (0,0) {zdrojový\\program};
            \node<2->[box] (lex) at (3.25,0) {lexikální\\analyzátor};
            \node<4->[box] (tok) at (6.50,0) {posloupnost\\tokenů};

            \draw<3->[arrow] (src) -- (lex);
            \draw<4->[arrow] (lex) -- (tok);

            \node<6->[font=\scriptsize, align=center] at (0,-0.9)
                {\texttt{position = speed * 60;}};

            \node<6->[font=\scriptsize, align=center] at (6.50,-0.9)
                {\texttt{id = id * num ;}};

        \end{tikzpicture}

    \end{frame}

    \subsection{Lexémy a tokeny}

    \begin{frame}{Lexém a token}

        \begin{itemize}
            \visible<2->{\item \textbf{Lexém} je konkrétní část zdrojového textu.}
            \visible<3->{\item \textbf{Token} je informace, kterou dostane další část překladače.}
            \visible<4->{\item Token často zapisujeme jako dvojici}
        \end{itemize}

        \visible<4->{
            \[
                \langle \text{typ tokenu}, \text{hodnota} \rangle
            \]
        }

        \vspace{-0.1cm}

        \centering
        \begin{tikzpicture}
            \node<5->[
                draw,
                thick,
                rounded corners,
                fill=blue!20!white,
                inner xsep=8pt,
                inner ysep=6pt
            ] {\texttt{position = speed * 60;}};
        \end{tikzpicture}

        \vspace{0.25cm}

        \visible<5->{
            \begin{center}
                \scriptsize
                \renewcommand{\arraystretch}{1.25}
                \begin{tabular}{|c|c|c|}
                    \hline
                    \textbf{lexém} & \textbf{typ tokenu} & \textbf{token} \\
                    \hline
                    \visible<6->{\texttt{position}} & \visible<6->{identifikátor} & \visible<6->{\texttt{<id, position>}} \\
                    \hline
                    \visible<7->{\texttt{=}} & \visible<7->{operátor přiřazení} & \visible<7->{\texttt{<=>}} \\
                    \hline
                    \visible<8->{\texttt{speed}} & \visible<8->{identifikátor} & \visible<8->{\texttt{<id, speed>}} \\
                    \hline
                    \visible<9->{\texttt{*}} & \visible<9->{operátor} & \visible<9->{\texttt{<*>}} \\
                    \hline
                    \visible<10->{\texttt{60}} & \visible<10->{číslo} & \visible<10->{\texttt{<num, 60>}} \\
                    \hline
                    \visible<11->{\texttt{;}} & \visible<11->{oddělovač} & \visible<11->{\texttt{<;>}} \\
                    \hline
                \end{tabular}
            \end{center}
        }

    \end{frame}

    \begin{frame}{Od regulárního výrazu k~lexeru}

        \begin{itemize}
            \visible<2->{\item Každý typ tokenu popíšeme jako \textbf{regulární jazyk}.}
            \visible<3->{\item Pro tyto jazyky lze sestrojit \textbf{konečné automaty}.}
            \visible<4->{\item Výsledný automat pak čte zdrojový program znak po znaku.}
        \end{itemize}

        \vspace{0.05cm}

        \centering
        \begin{tikzpicture}[
            box/.style={
                draw,
                rounded corners,
                thick,
                align=center,
                font=\scriptsize,
                minimum height=0.72cm
            },
            smallbox/.style={
                draw,
                rounded corners,
                thick,
                align=center,
                font=\scriptsize,
                minimum width=1.45cm,
                minimum height=0.62cm
            },
            arrow/.style={->, thick}
        ]

            % horní řada: regulární výrazy tokenů
            \node<2->[box, minimum width=2.45cm] (re1) at (0,0)
                {identifikátor\\\texttt{[a-zA-Z][a-zA-Z0-9]*}};

            \node<2->[box, minimum width=1.95cm] (re2) at (3.05,0)
                {číslo\\\texttt{[0-9]+}};

            \node<2->[box, minimum width=2.15cm] (re3) at (5.80,0)
                {operátor\\\texttt{= \;|\; + \;|\; *}};

            % mezikrok: společný regulární popis
            \node<3->[box, minimum width=4.8cm] (combined) at (2.90,-1.45)
                {společný regulární popis\\
                \texttt{id \;|\; num \;|\; op}};

            \draw<3->[arrow] (re1.south) -- (combined.north west);
            \draw<3->[arrow] (re2.south) -- (combined.north);
            \draw<3->[arrow] (re3.south) -- (combined.north east);

            % konstrukce automatů
            \node<4->[smallbox] (nfa) at (0.60,-2.95) {NKA};
            \node<5->[smallbox] (dfa) at (3.50,-2.95) {DKA};
            \node<6->[box, minimum width=2.05cm] (lexer) at (6.00,-2.95)
                {lexikální\\analyzátor};

            \draw<4->[arrow] (combined.south) -- ++(0,-0.225) -| (nfa.north);
            \draw<5->[arrow] (nfa) -- node[above, font=\scriptsize, yshift=1mm] {determin.} (dfa);
            \draw<6->[arrow] (dfa) -- (lexer);

            % spodní ukázka
            \node<7->[box, minimum width=2.4cm] (src) at (0.65,-4.15)
                {\texttt{position = 60;}};

            \node<7->[box, minimum width=3.3cm] (tok) at (5.35,-4.15)
                {\texttt{<id, position>}\\
                \texttt{<=> <num, 60> <;>}};

            \draw<7->[arrow] (src) -- node[above, font=\scriptsize, yshift=1mm] {čtení znaků} (tok);

        \end{tikzpicture}

    \end{frame}

    \begin{frame}{Obecná konstrukce lexeru}

        \small
        \begin{itemize}
            \visible<2->{\item Pro každý typ tokenu sestrojíme automat \(T_1,\dots,T_n\).}
            \visible<3->{\item Přidáme společný počáteční stav \(s\) a \(\lambda\)-přechody;
            mezery a komentáře pouze přeskočíme.}
        \end{itemize}

        \vspace{-0.05cm}

        \centering
        \begin{tikzpicture}[scale=0.82, every node/.style={font=\scriptsize}]

            %-------------------------------------------------
            % společný počáteční stav
            %-------------------------------------------------
            \initialstate<2->{S}{(-1,0)}{$s$}{(-2.0,0.0)}{west}
            \transitionarrow<3->[loop above]{S}{above}{\scriptsize mezera, tab, EOL}{S}

            %-------------------------------------------------
            % větev pro komentář
            %-------------------------------------------------
            % \draw<3->[green!60!black, thick, rounded corners]
            %     (-1.15,-2.55) rectangle (1.20,-1.75);

            \state<3->{C1}{(-1.0,-1.80)}{$c_1$}
            \acceptingstate<3->{C2}{(-1.0,-3.60)}{$c_2$}

            \transitionarrow<3->{S}{left}{\scriptsize zač. komentáře}{C1}
            \transitionarrow<3->[loop right]{C1}{below}{\scriptsize obsah}{C1}
            \transitionarrow<3->{C1}{right}{\scriptsize konec}{C2}
            \transitionarrow<3->[bend left=70]{C2}{left}{\scriptsize přeskočit}{S}

            %-------------------------------------------------
            % T1
            %-------------------------------------------------
            \draw<2->[cyan!70!black, thick, rounded corners]
                (3.00,0.8) rectangle (10.20,-0.80);

            \node<2->[text=cyan!50!black]
                at (10.60,0.0) {$T_1$};

            \state<2->{Q1}{(4.10,0)}{$q_1$}
            \node<2-> (T1dots) at (6.60,0) {\(\cdots\)};
            \acceptingstate<2->{F1}{(9.10,0)}{$f_1$}

            \transitionarrow<2->{S}{above}{\scriptsize \(\lambda\)}{Q1}
            \transitionarrow<2->{Q1}{above}{}{T1dots}
            \transitionarrow<2->{T1dots}{above}{}{F1}

            %-------------------------------------------------
            % T2
            %-------------------------------------------------
            \draw<2->[cyan!70!black, thick, rounded corners]
                (3.00,-1.25) rectangle (10.20,-2.85);

            \node<2->[text=cyan!50!black]
                at (10.60,-2.05) {$T_2$};

            \state<2->{Q2}{(4.10,-2.05)}{$q_2$}
            \node<2-> (T2dots) at (6.60,-2.05) {\(\cdots\)};
            \acceptingstate<2->{F2}{(9.10,-2.05)}{$f_2$}

            \transitionarrow<2->{S}{above}{\scriptsize \(\lambda\)}{Q2}
            \transitionarrow<2->{Q2}{above}{}{T2dots}
            \transitionarrow<2->{T2dots}{above}{}{F2}

            %-------------------------------------------------
            % Tn
            %-------------------------------------------------
            \draw<2->[cyan!70!black, thick, rounded corners]
                (3.00,-3.40) rectangle (10.20,-5.00);

            \node<2->[text=cyan!50!black]
                at (10.60,-4.20) {$T_n$};

            \state<2->{Qn}{(4.10,-4.25)}{$q_n$}
            \node<2-> (Tndots) at (6.60,-4.25) {\(\cdots\)};
            \acceptingstate<2->{Fn}{(9.10,-4.25)}{$f_n$}

            \transitionarrow<2->{S}{above}{\scriptsize \(\lambda\)}{Qn}
            \transitionarrow<2->{Qn}{above}{}{Tndots}
            \transitionarrow<2->{Tndots}{above}{}{Fn}

            %-------------------------------------------------
            % naznačení dalších větví
            %-------------------------------------------------
            \node<2->[font=\Large] at (4.05,-3.15) {\(\vdots\)};

        \end{tikzpicture}

    \end{frame}

    \begin{frame}{Jak lexer vybere správný token?}

        \begin{itemize}
            \visible<2->{\item Lexer čte vstup zleva doprava a hledá nejdelší prefix,
            který tvoří platný lexém.}
            \visible<3->{\item Pokud může vytvořit delší lexém, pokračuje ve čtení.}
            \visible<4->{\item Pokud se na stejný lexém hodí více pravidel,
            rozhodne \textbf{priorita pravidel}.}
        \end{itemize}

        \vspace{0.10cm}

        \centering
        \renewcommand{\arraystretch}{1.25}
        \visible<5->{
            \begin{tabular}{|c|c|c|}
                \hline
                \textbf{vstup} & \textbf{během čtení lze přijmout} & \textbf{výsledek lexeru} \\
                \hline
                \visible<5->{\texttt{==}} &
                \visible<5->{\texttt{=} i \texttt{==}} &
                \visible<5->{\texttt{==}} \\
                \hline
                \visible<6->{\texttt{abc123}} &
                \visible<6->{\texttt{a}, \texttt{ab}, \dots, \texttt{abc123}} &
                \visible<6->{\texttt{<id, abc123>}} \\
                \hline
                \visible<7->{\texttt{then}} &
                \visible<7->{klíčové slovo i identifikátor} &
                \visible<7->{klíčové slovo} \\
                \hline
            \end{tabular}
        }

    \end{frame}

    \section{Syntaktická analýza}

    \titleframe{Syntaktická analýza}

    \subsection{Co je to parser?}

    \begin{frame}{Syntaktická analýza}

        \begin{itemize}
            \visible<2->{\item \textbf{Syntaktický analyzátor} neboli \textbf{parser}
            kontroluje, zda posloupnost tokenů tvoří syntakticky správný program.}
            \visible<3->{\item Vstupem parseru jsou \textbf{tokeny} od lexikálního analyzátoru (lexeru).}
            \visible<4->{\item Parser se řídí \textbf{gramatikou} daného programovacího jazyka.}
            \visible<5->{\item Výsledkem je typicky informace}
            \visible<6->{
                \begin{itemize}
                    \item že je vstup syntakticky správný,
                    \item nebo že obsahuje syntaktickou chybu,
                    \item případně také derivační strom.
                \end{itemize}
            }
        \end{itemize}

        \vspace{0.1cm}

        \centering
        \begin{tikzpicture}[
            box/.style={
                draw,
                rounded corners,
                thick,
                align=center,
                minimum width=2.2cm,
                minimum height=0.8cm,
                font=\scriptsize
            },
            arrow/.style={->, thick}
        ]
            \node<3->[box] (lex) at (0,0) {tokeny};
            \node<2->[box] (parser) at (3.2,0) {syntaktický\\analyzátor};
            \node<5->[box] (out) at (6.4,0) {správný vstup\\/ chyba / strom};

            \draw<3->[arrow] (lex) -- (parser);
            \draw<5->[arrow] (parser) -- (out);
        \end{tikzpicture}

    \end{frame}

    \subsection{Jak parser pracuje s tokeny?}

    \begin{frame}{Jak parser pracuje s tokeny?}

        \begin{itemize}
            \visible<2->{\item Lexer už rozdělil zdrojový text na tokeny.}
            \visible<3->{\item Parser už nepracuje s jednotlivými znaky,
            ale s \textbf{posloupností tokenů}.}
            \visible<4->{\item Snaží se zjistit, zda tuto posloupnost lze odvodit
            z počátečního symbolu gramatiky.}
        \end{itemize}

        \vspace{0.1cm}

        \centering
        \begin{tikzpicture}[
            tok/.style={
                draw,
                rounded corners,
                thick,
                align=center,
                minimum width=1.1cm,
                minimum height=0.65cm,
                font=\scriptsize
            },
            box/.style={
                draw,
                rounded corners,
                thick,
                align=center,
                minimum width=2.1cm,
                minimum height=0.8cm,
                font=\scriptsize
            },
            arrow/.style={->, thick}
        ]
            \node<2->[tok] (t1) at (0,0) {\texttt{<id>}};
            \node<2->[tok] (t2) at (1.5,0) {\texttt{<+>}};
            \node<2->[tok] (t3) at (3.0,0) {\texttt{<id>}};
            \node<2->[tok] (t4) at (4.5,0) {\texttt{<*>}};
            \node<2->[tok] (t5) at (6.0,0) {\texttt{<num>}};

            \node<3->[box] (parser) at (3.0,-1.2) {parser};

            \draw<3->[arrow] (t1.south) -- ++(0,-0.35) -- (parser.north west);
            \draw<3->[arrow] (t2.south) -- ++(0,-0.30) -- (parser.north west);
            \draw<3->[arrow] (t3.south) -- (parser.north);
            \draw<3->[arrow] (t4.south) -- ++(0,-0.30) -- (parser.north east);
            \draw<3->[arrow] (t5.south) -- ++(0,-0.35) -- (parser.north east);

            \node<4->[font=\scriptsize, align=center] at (3.0,-2.1)
                {např. ověří, že jde o aritmetický výraz};
        \end{tikzpicture}

    \end{frame}

    \subsection{Bezkontextové gramatiky}

    \begin{frame}{Bezkontextové gramatiky}

        Syntaxi programovacího jazyka obvykle popisujeme pomocí \textbf{bezkontextové gramatiky}.
        \begin{defblock}[Bezkontextová gramatika]
            Gramatika $G=(V,T,P,S)$ se nazývá \markdarkred{bezkontextová}, pokud $P$ obsahuje pouze pravidla tvaru \(A \to \alpha\), kde \(A \in V\) a~\(\alpha \in (V \cup T)^*\).
        \end{defblock}

    \end{frame}

    \begin{frame}{Příklad bezkontextové gramatiky}{Jazyk $a^n b^n$}

        \visible<2->{
            Uvažujme gramatiku $G=(V,T,P,S)$, kde $V=\set{S},T=\set{a,b}$ a pravidla jsou
            \[
                S \to aSb \mid \lambda.
            \]
        }
        \visible<3->{Pak $G$ generuje jazyk $L(G)=\set{a^n b^n \mid n \geqslant 0}$.}

        \[
            \visible<4->{S}
            \visible<5->{\Rightarrow aSb}
            \visible<6->{\Rightarrow aaSbb}
            \visible<7->{\Rightarrow aa\lambda bb}
            \visible<8->{= aabb}
        \]

    \end{frame}

    \subsection{Derivační strom}

    \begin{frame}{Derivační strom}

        \begin{itemize}
            \visible<2->{\item Derivační strom názorně ukazuje,
            jak byl řetězec odvozen z gramatiky.}
            \visible<3->{\item Vnitřní vrcholy odpovídají neterminálům,}
            \visible<4->{\item listy odpovídají terminálům výsledného řetězce.}
        \end{itemize}

        \vspace{0.15cm}

        \centering
        \begin{tikzpicture}[
            node distance=0.8cm,
            every node/.style={font=\scriptsize},
            nt/.style={draw, circle, thick, minimum size=7mm},
            t/.style={draw, rounded corners, thick, minimum width=6mm, minimum height=5mm}
        ]
            \node<5->[nt] (S) at (0,0) {$S$};
            \node<6->[t] (a1) at (-1.4,-1.1) {$a$};
            \node<6->[nt] (S2) at (0,-1.1) {$S$};
            \node<6->[t] (b1) at (1.4,-1.1) {$b$};

            \node<7->[t] (a2) at (-0.9,-2.2) {$a$};
            \node<7->[nt] (S3) at (0,-2.2) {$S$};
            \node<7->[t] (b2) at (0.9,-2.2) {$b$};

            \node<8->[t] (eps) at (0,-3.2) {$\lambda$};

            \draw<6->[thick] (S) -- (a1);
            \draw<6->[thick] (S) -- (S2);
            \draw<6->[thick] (S) -- (b1);

            \draw<7->[thick] (S2) -- (a2);
            \draw<7->[thick] (S2) -- (S3);
            \draw<7->[thick] (S2) -- (b2);

            \draw<8->[thick] (S3) -- (eps);
        \end{tikzpicture}

        \vspace{0.05cm}

        \visible<9->{
            \begin{center}
                \scriptsize
                Derivační strom slova $aabb$ podle předchozí gramatiky $G$.
            \end{center}
        }

    \end{frame}

    \subsection{Aritmetické výrazy}

    \begin{frame}{Gramatika aritmetických výrazů}

        \begin{columns}[T,onlytextwidth]
            \begin{column}{0.48\textwidth}
                \visible<2->{Definujeme gramatiku generující aritmetické výrazy s~operátory \(+,-,*,/\) a~závorkami.
                    \[
                    \begin{aligned}
                        E &\to E + T \mid E - T \mid T \\
                        T &\to T * F \mid T / F \mid F \\
                        F &\to (E) \mid \mathtt{id} \mid \mathtt{num}
                    \end{aligned}
                    \]
                }
            \end{column}

            \begin{column}{0.48\textwidth}
                \begin{itemize}
                    \visible<3->{\item \(E\) = výraz}
                    \visible<4->{\item \(T\) = term}
                    \visible<5->{\item \(F\) = faktor}
                    \visible<6->{\item gramatika zachycuje prioritu operátorů}
                    \visible<7->{\item násobení a dělení mají vyšší prioritu než sčítání a odčítání}
                \end{itemize}
            \end{column}
        \end{columns}

        \vspace{0.4cm}

        \onslide<8->
        \textbf{Příklady správných výrazů:} \texttt{a + b}, \texttt{a * (b + 3)}, \texttt{(a + b) * c}

    \end{frame}

    \begin{frame}{Příklad derivačního stromu}{Výraz \texttt{id + id * id}}

        \begin{columns}[T,onlytextwidth]
            \begin{column}{0.3\textwidth}
                \pause
                \[
                \begin{aligned}
                    E &\to E + T \mid E - T \mid T \\
                    T &\to T * F \mid T / F \mid F \\
                    F &\to (E) \mid \mathtt{id} \mid \mathtt{num}
                \end{aligned}
                \]
            \end{column}

            \begin{column}{0.7\textwidth}
                \centering
                \begin{tikzpicture}[
                    every node/.style={font=\scriptsize},
                    nt/.style={draw, circle, thick, minimum size=7mm},
                    t/.style={draw, rounded corners, thick, minimum width=6mm, minimum height=5mm}
                ]
                    \node<2->[nt] (E0) at (0,0) {$E$};

                    \node<3->[nt] (E1) at (-2.2,-1.1) {$E$};
                    \node<3->[t]  (plus) at (0,-1.1) {$+$};
                    \node<3->[nt] (T1) at (2.2,-1.1) {$T$};

                    \draw<3->[thick] (E0) -- (E1);
                    \draw<3->[thick] (E0) -- (plus);
                    \draw<3->[thick] (E0) -- (T1);

                    \node<4->[nt] (Tleft) at (-2.2,-2.2) {$T$};
                    \node<4->[nt] (T2) at (1.2,-2.2) {$T$};
                    \node<4->[t]  (times) at (2.2,-2.2) {$*$};
                    \node<4->[nt] (F3) at (3.2,-2.2) {$F$};

                    \draw<4->[thick] (E1) -- (Tleft);
                    \draw<4->[thick] (T1) -- (T2);
                    \draw<4->[thick] (T1) -- (times);
                    \draw<4->[thick] (T1) -- (F3);

                    \node<5->[nt] (Fleft) at (-2.2,-3.3) {$F$};
                    \node<5->[nt] (F2) at (1.2,-3.3) {$F$};

                    \draw<5->[thick] (Tleft) -- (Fleft);
                    \draw<5->[thick] (T2) -- (F2);

                    \node<6->[t] (a1) at (-2.2,-4.2) {$\mathtt{id}$};
                    \node<6->[t] (a2) at (1.2,-4.2) {$\mathtt{id}$};
                    \node<6->[t] (a3) at (3.2,-4.2) {$\mathtt{id}$};

                    \draw<6->[thick] (Fleft) -- (a1);
                    \draw<6->[thick] (F2) -- (a2);
                    \draw<6->[thick] (F3) -- (a3);
                \end{tikzpicture}
            \end{column}
        \end{columns}

        \vspace{0.2cm}

        \onslide<7->
        \begin{timportantframe}[0.6\textwidth]
            \begin{center}
                Strom odpovídá slovu $\mathtt{id} + \mathtt{id} * \mathtt{id}$.
            \end{center}
        \end{timportantframe}

    \end{frame}

    \begin{frame}{Příklad derivačního stromu}{Výraz \texttt{(id + num) * id}}

        \begin{columns}[T,onlytextwidth]
            \begin{column}{0.4\textwidth}
                \centering
                \begin{tikzpicture}[
                    scale=0.78,
                    every node/.style={font=\scriptsize},
                    nt/.style={draw, circle, thick, minimum size=7mm},
                    t/.style={draw, rounded corners, thick, minimum width=7mm, minimum height=5mm}
                ]

                    \node<2->[nt] (E0) at (0,0) {$E$};

                    \node<3->[nt] (T0) at (0,-1.0) {$T$};
                    \draw<3->[thick] (E0) -- (T0);

                    \node<4->[nt] (T1) at (-1.3,-2.0) {$T$};
                    \node<4->[t]  (times) at (0,-2.0) {$*$};
                    \node<4->[nt] (F3) at (1.3,-2.0) {$F$};

                    \draw<4->[thick] (T0) -- (T1);
                    \draw<4->[thick] (T0) -- (times);
                    \draw<4->[thick] (T0) -- (F3);

                    \node<5->[nt] (F1) at (-1.3,-3.0) {$F$};
                    \node<5->[t]  (id3) at (1.3,-3.0) {$\mathtt{id}$};

                    \draw<5->[thick] (T1) -- (F1);
                    \draw<5->[thick] (F3) -- (id3);

                    \node<6->[t]  (lpar) at (-2.7,-4.0) {$($};
                    \node<6->[nt] (E1) at (-1.3,-4.0) {$E$};
                    \node<6->[t]  (rpar) at (0.1,-4.0) {$)$};

                    \draw<6->[thick] (F1) -- (lpar);
                    \draw<6->[thick] (F1) -- (E1);
                    \draw<6->[thick] (F1) -- (rpar);

                    \node<7->[nt] (E2) at (-2.3,-5.0) {$E$};
                    \node<7->[t]  (plus) at (-1.3,-5.0) {$+$};
                    \node<7->[nt] (T2) at (-0.3,-5.0) {$T$};

                    \draw<7->[thick] (E1) -- (E2);
                    \draw<7->[thick] (E1) -- (plus);
                    \draw<7->[thick] (E1) -- (T2);

                    \node<8->[nt] (T3) at (-2.3,-6.0) {$T$};
                    \node<8->[nt] (F2) at (-0.3,-6.0) {$F$};

                    \draw<8->[thick] (E2) -- (T3);
                    \draw<8->[thick] (T2) -- (F2);

                    \node<9->[nt] (F4) at (-2.3,-7.0) {$F$};
                    \node<9->[t] (num1) at (-0.3,-7.0) {$\mathtt{num}$};

                    \draw<9->[thick] (T3) -- (F4);
                    \draw<9->[thick] (F2) -- (num1);

                    \node<10->[t] (id1) at (-2.3,-7.9) {$\mathtt{id}$};
                    \draw<10->[thick] (F4) -- (id1);

                \end{tikzpicture}
            \end{column}

            \begin{column}{0.58\textwidth}
                \begin{align*}
                    \small
                    E
                    &\visible<3->{\Rightarrow T}
                    \visible<4->{\Rightarrow T * F}
                    \visible<5->{\Rightarrow F * F} \\
                    &\visible<6->{\Rightarrow (E) * F}
                    \visible<7->{\Rightarrow (E + T) * F} \\
                    &\visible<8->{\Rightarrow (T + F) * F}
                    \visible<9->{\Rightarrow (F + \mathtt{num}) * F} \\
                    &\visible<10->{\Rightarrow (\mathtt{id} + \mathtt{num}) * F}
                    \visible<11->{\Rightarrow (\mathtt{id} + \mathtt{num}) * \mathtt{id}}
                \end{align*}

                \onslide<12->
                \begin{timportantframe}
                    Strom odpovídá například zdrojovému výrazu \texttt{(a + 3) * b},
                    který lexer převede na
                    \((\mathtt{id} + \mathtt{num}) * \mathtt{id}\).
                \end{timportantframe}
            \end{column}
        \end{columns}

    \end{frame}

    \begin{frame}{Gramatika jednoduchého příkazu v C\#}
        \centering
        \begin{tikzpicture}
            \node[
                draw,
                thick,
                rounded corners,
                fill=blue!20!white,
                inner xsep=8pt,
                inner ysep=6pt
            ] {\texttt{float posX = currPosX + stepsX * directionX * speedX;}};
        \end{tikzpicture}
        \begin{columns}[T,onlytextwidth]
            \begin{column}{0.52\textwidth}
                \[
                \begin{aligned}
                    S &\to \mathit{Type}\ \mathtt{id}\ \mathtt{=}\ E\ \mathtt{;} \\
                    \mathit{Type} &\to \mathtt{int} \mid \mathtt{double} \mid \mathtt{float} \\
                    E &\to E\ \mathtt{+}\ T \mid E\ \mathtt{-}\ T \mid T \\
                    T &\to T\ \mathtt{*}\ F \mid T\ \mathtt{/}\ F \mid F \\
                    F &\to \mathtt{id} \mid \mathtt{num} \mid \mathtt{(}E\mathtt{)}
                \end{aligned}
                \]
            \end{column}

            \begin{column}{0.46\textwidth}
                \textbf{Tokeny pro parser:}
                \[
                \begin{aligned}
                    &\mathtt{float}\ \mathtt{id}\ \mathtt{=}\\
                    &\mathtt{id}\ \mathtt{+}\ \mathtt{id}\ \mathtt{*}\\
                    &\mathtt{id}\ \mathtt{*}\ \mathtt{id}\ \mathtt{;}
                \end{aligned}
                \]
            \end{column}
        \end{columns}

    \end{frame}

    \begin{frame}{Derivační strom}{Příkaz \texttt{float posX = currPosX + stepsX * directionX * speedX;}}

        \centering
        \begin{tikzpicture}[
            scale=0.85,
            transform shape,
            every node/.style={font=\scriptsize},
            nt/.style={draw, circle, thick, minimum size=8mm},
            t/.style={draw, rounded corners, thick, minimum width=8mm, minimum height=6mm}
        ]

            % vrcholy
            \node[nt] (S) at (0,0) {$S$};

            \node[nt] (Type) at (-5.6,-1.3) {$Type$};
            \node[t]  (id0)  at (-3.2,-1.3) {$\mathtt{id}$};
            \node[t]  (eq)   at (-1.5,-1.3) {$=$};
            \node[nt] (E0)   at (1.0,-1.3) {$E$};
            \node[t]  (semi) at (5.8,-1.3) {$;$};

            \node[t]  (float)   at (-5.6,-2.5) {$\mathtt{float}$};

            \node[nt] (E1)    at (-0.8,-2.5) {$E$};
            \node[t]  (plus)  at (1.0,-2.5) {$+$};
            \node[nt] (T0)    at (3.0,-2.5) {$T$};

            \node[nt] (Tleft)  at (-0.8,-3.7) {$T$};

            \node[nt] (T1)     at (1.6,-3.7) {$T$};
            \node[t]  (mul2)   at (3.0,-3.7) {$*$};
            \node[nt] (Fspeed) at (4.4,-3.7) {$F$};

            \node[nt] (Fleft)   at (-0.8,-4.9) {$F$};

            \node[nt] (T2)      at (0.2,-4.9) {$T$};
            \node[t]  (mul1)    at (1.6,-4.9) {$*$};
            \node[nt] (Fdir)    at (3.0,-4.9) {$F$};

            \node[t]  (idSpeed) at (4.4,-4.9) {$\mathtt{id}$};

            \node[t]  (idCurr)  at (-0.8,-6.1) {$\mathtt{id}$};

            \node[nt] (Fsteps)  at (0.2,-6.1) {$F$};
            \node[t]  (idDir)   at (3.0,-6.1) {$\mathtt{id}$};

            \node[t]  (idSteps) at (0.2,-7.3) {$\mathtt{id}$};

            % hrany
            \draw[thick] (S) -- (Type);
            \draw[thick] (S) -- (id0);
            \draw[thick] (S) -- (eq);
            \draw[thick] (S) -- (E0);
            \draw[thick] (S) -- (semi);

            \draw[thick] (Type) -- (float);

            \draw[thick] (E0) -- (E1);
            \draw[thick] (E0) -- (plus);
            \draw[thick] (E0) -- (T0);

            \draw[thick] (E1) -- (Tleft);
            \draw[thick] (Tleft) -- (Fleft);
            \draw[thick] (Fleft) -- (idCurr);

            \draw[thick] (T0) -- (T1);
            \draw[thick] (T0) -- (mul2);
            \draw[thick] (T0) -- (Fspeed);

            \draw[thick] (T1) -- (T2);
            \draw[thick] (T1) -- (mul1);
            \draw[thick] (T1) -- (Fdir);

            \draw[thick] (T2) -- (Fsteps);
            \draw[thick] (Fsteps) -- (idSteps);

            \draw[thick] (Fdir) -- (idDir);
            \draw[thick] (Fspeed) -- (idSpeed);

        \end{tikzpicture}

    \end{frame}

    \subsection{Analýza shora dolů}

    \begin{frame}{Analýza shora dolů}

        \begin{itemize}
            \visible<2->{\item Parser začíná počátečním symbolem gramatiky.}
            \visible<3->{\item Postupně rozvíjí neterminály tak,
            aby odvodil vstupní posloupnost tokenů.}
            \visible<4->{\item Strom tedy vzniká \textbf{od kořene k listům}.}
        \end{itemize}

        \vspace{0.2cm}

        \centering
        \begin{tikzpicture}[
            box/.style={
                draw,
                rounded corners,
                thick,
                minimum width=1.4cm,
                minimum height=0.6cm,
                align=center,
                font=\scriptsize
            },
            arrow/.style={->, thick}
        ]
            \node<5->[box] (s1) at (0,0) {$E$};
            \node<6->[box] (s2) at (2.0,0) {$E + T$};
            \node<7->[box] (s3) at (4.4,0) {$T + T$};
            \node<8->[box] (s4) at (6.8,0) {$F + F$};
            \node<9->[box] (s5) at (9.2,0) {$a + a$};

            \draw<6->[arrow] (s1) -- (s2);
            \draw<7->[arrow] (s2) -- (s3);
            \draw<8->[arrow] (s3) -- (s4);
            \draw<9->[arrow] (s4) -- (s5);
        \end{tikzpicture}

        \vspace{0.2cm}

        \onslide<10->
        \begin{tinfoframe}
            Shora dolů tedy "předpovídáme", jaká pravidla bude třeba použít.
        \end{tinfoframe}

    \end{frame}

    \subsection{Analýza zdola nahorů}

    \begin{frame}{Analýza zdola nahoru}

        \begin{itemize}
            \visible<2->{\item Parser začíná vstupní posloupností tokenů.}
            \visible<3->{\item Postupně hledá části vstupu, které odpovídají pravým stranám pravidel,
            a nahrazuje je levými stranami.}
            \visible<4->{\item Strom tedy vzniká \textbf{od listů ke kořeni}.}
        \end{itemize}

        \vspace{0.3cm}

        \centering
        \begin{tikzpicture}[
            box/.style={
                draw,
                rounded corners,
                thick,
                minimum width=1.4cm,
                minimum height=0.6cm,
                align=center,
                font=\scriptsize
            },
            arrow/.style={->, thick}
        ]
            \node<5->[box] (s1) at (0,0) {$a + a$};
            \node<6->[box] (s2) at (2.2,0) {$F + F$};
            \node<7->[box] (s3) at (4.6,0) {$T + T$};
            \node<8->[box] (s4) at (7.0,0) {$E + T$};
            \node<9->[box] (s5) at (9.4,0) {$E$};

            \draw<6->[arrow] (s1) -- (s2);
            \draw<7->[arrow] (s2) -- (s3);
            \draw<8->[arrow] (s3) -- (s4);
            \draw<9->[arrow] (s4) -- (s5);
        \end{tikzpicture}

        \vspace{0.3cm}

        \onslide<10->
        \begin{tinfoframe}
            Zdola nahoru tedy postupně \textbf{redukujeme} vstup až na počáteční symbol.
        \end{tinfoframe}

    \end{frame}

    \section{Sémantická analýza}

    \titleframe{Sémantická analýza}

    \begin{frame}{Sémantická analýza}

        \begin{itemize}
            \visible<2->{\item Parser ověří, zda má program správnou \textbf{syntaxi}.}
            \visible<3->{\item \textbf{Sémantická analýza} ověřuje, zda má program správný \textbf{význam}.}
            \visible<4->{\item Kontroluje například deklarace proměnných, typy výrazů a platnost operací.}
        \end{itemize}

        \vspace{0.3cm}

        \visible<5->{\centering
        \begin{tikzpicture}[
            box/.style={
                draw,
                rounded corners,
                thick,
                align=center,
                font=\scriptsize,
                minimum width=2.1cm,
                minimum height=0.75cm
            },
            arrow/.style={->, thick}
        ]

            \node[box] (parser) at (0,0) {parser};
            \node[box] (tree) at (2.7,0) {syntaktický\\strom};
            \node[box] (sem) at (5.4,0) {sémantická\\analýza};
            \node[box] (out) at (8.1,0) {anotovaný strom\\nebo chyba};

            \draw[arrow] (parser) -- (tree);
            \draw[arrow] (tree) -- (sem);
            \draw[arrow] (sem) -- (out);

        \end{tikzpicture}}

    \end{frame}

    \begin{frame}{Syntaxe vs. sémantika}

        \begin{itemize}
            \visible<2->{\item \textbf{Syntaktická chyba}: program porušuje pravidla zápisu.}
            \visible<3->{\item \textbf{Sémantická chyba}: program je zapsán správně, ale nedává smysl.}
        \end{itemize}

        \vspace{0.3cm}

        \centering\small
        \renewcommand{\arraystretch}{1.2}
        \begin{tabular}{|c|c|c|}
            \hline
            \textbf{příkaz} & \textbf{syntaxe} & \textbf{sémantika} \\
            \hline
            \visible<4->{\texttt{int x = 5;}} &
            \visible<4->{správně} &
            \visible<4->{správně} \\
            \hline
            \visible<5->{\texttt{int x = ;}} &
            \visible<5->{chyba} &
            \visible<5->{--} \\
            \hline
            \visible<6->{\texttt{string s = 5;}} &
            \visible<6->{správně} &
            \visible<6->{chyba typu} \\
            \hline
            \visible<7->{\texttt{x = y + 1;}} &
            \visible<7->{správně} &
            \visible<7->{chyba, pokud \(x\) nebo \(y\) nejsou deklarované} \\
            \hline
        \end{tabular}

        \vspace{0.15cm}

        \onslide<8->
        \begin{timportantframe}
            Syntaxe říká, \textbf{jak má být program zapsán}.  
            Sémantika říká, \textbf{zda tento zápis dává smysl}.
        \end{timportantframe}

    \end{frame}

    \begin{frame}{Tabulka symbolů}

        \begin{itemize}
            \visible<2->{\item Sémantická analýza pracuje s \textbf{tabulkou symbolů}.}
            \visible<3->{\item Ukládá si do ní informace o proměnných, funkcích a jejich typech.}
        \end{itemize}

        \vspace{0.1cm}

        \visible<4->{
            \centering
            \begin{tikzpicture}
                \node[
                    draw,
                    thick,
                    rounded corners,
                    fill=blue!20!white,
                    inner xsep=8pt,
                    inner ysep=6pt
                ] {\texttt{float posX = currPosX + stepsX * directionX * speedX;}};
            \end{tikzpicture}
        }

        \vspace{0.3cm}

        \centering
        \visible<5->{
            \small
            \renewcommand{\arraystretch}{1.2}
            \begin{tabular}{|c|c|c|}
                \hline
                \textbf{identifikátor} & \textbf{typ} & \textbf{poznámka} \\
                \hline
                \texttt{posX} & \texttt{int} & deklarovaná proměnná \\
                \hline
                \texttt{currPosX} & \texttt{int} & musí být deklarována dříve \\
                \hline
                \texttt{stepsX} & \texttt{int} & musí být deklarována dříve \\
                \hline
                \texttt{directionX} & \texttt{int} & musí být deklarována dříve \\
                \hline
                \texttt{speedX} & \texttt{int} & musí být deklarována dříve \\
                \hline
            \end{tabular}
        }

    \end{frame}

    \begin{frame}{Typová kontrola výrazu}

        \begin{itemize}
            \visible<2->{\item Sémantická analýza přiřazuje částem výrazu typy.}
            \visible<3->{\item Operace je validní jen tehdy, pokud odpovídají typy operandů.}
        \end{itemize}

        \vspace{0.1cm}

        \centering
        \begin{tikzpicture}[
            every node/.style={font=\scriptsize},
            box/.style={
                draw,
                rounded corners,
                thick,
                align=center,
                minimum width=1.8cm,
                minimum height=0.55cm
            },
            arrow/.style={->, thick}
        ]

            \node<4->[box] (a) at (0,0) {\texttt{stepsX}\\\texttt{int}};
            \node<4->[box] (b) at (2.2,0) {\texttt{directionX}\\\texttt{int}};
            \node<5->[box] (mul1) at (1.1,-1.0) {\texttt{*}\\výsledek \texttt{int}};

            \draw<5->[arrow] (a) -- (mul1);
            \draw<5->[arrow] (b) -- (mul1);

            \node<6->[box] (c) at (4.2,0) {\texttt{speedX}\\\texttt{float}};
            \node<6->[box] (mul2) at (2.7,-2.0) {\texttt{*}\\výsledek \texttt{float}};

            \draw<6->[arrow] (mul1) -- (mul2);
            \draw<6->[arrow] (c) -- (mul2);

            \node<6->[box] (d) at (-1.2,-2.0) {\texttt{currPosX}\\\texttt{float}};
            \node<7->[box] (plus) at (0.8,-3.0) {\texttt{+}\\výsledek \texttt{float}};

            \draw<7->[arrow] (d) -- (plus);
            \draw<7->[arrow] (mul2) -- (plus);

            \node<8->[box] (assign) at (0.8,-4.0) {\texttt{posX = ...}\\\texttt{float = float}};

            \draw<8->[arrow] (plus) -- (assign);

        \end{tikzpicture}

        \vspace{0.1cm}

        \onslide<8->
        \begin{timportantframe}
            \begin{center}
                Výraz má typ \texttt{float}, takže jej lze přiřadit do proměnné
                \texttt{posX} typu \texttt{float}.
            \end{center}
        \end{timportantframe}

    \end{frame}

    \section{Optimalizace}

    \titleframe{Optimalizace}

    \subsection{Optimalizace kódu}

    \begin{frame}{Optimalizace kódu}

        \begin{itemize}
            \visible<2->{\item Optimalizace navazuje na \textbf{přední část překladače}.}
            \visible<3->{\item Nejprve vznikne \textbf{mezikód}, nad kterým lze provádět úpravy.}
        \end{itemize}

        \vspace{0.1cm}

        \centering
        \begin{tikzpicture}[
            box/.style={
                draw,
                rounded corners,
                thick,
                align=center,
                font=\scriptsize,
                minimum width=1.75cm,
                minimum height=0.6cm
            },
            arrow/.style={->, thick}
        ]

            \node<2->[box, fill=red!35] (src) at (0,0) {zdrojový\\kód};
            \node<2->[box, fill=red!35] (lex) at (2.0,0) {lexikální\\analýza};
            \node<2->[box, fill=red!35] (syn) at (4.0,0) {syntaktická\\analýza};
            \node<2->[box, fill=red!35] (sem) at (6.0,0) {sémantická\\analýza};
            \node<3->[box, fill=yellow!45] (ir) at (6.0,-1.0) {mezikód};

            \draw<2->[arrow] (src) -- (lex);
            \draw<2->[arrow] (lex) -- (syn);
            \draw<2->[arrow] (syn) -- (sem);
            \draw<3->[arrow] (sem) -- (ir);

            \node<2->[font=\scriptsize\bfseries] at (3.0,0.65) {přední část překladače};

        \end{tikzpicture}

        \vspace{0.15cm}

        \begin{itemize}
            \visible<4->{\item \textbf{Optimalizace} upravuje mezikód tak,
            aby byl výsledný kód efektivnější.}
            \visible<5->{\item Může odstranit zbytečné výpočty, spočítat konstanty
            nebo nepočítat stejný výraz vícekrát.}
            \visible<6->{\item Nesmí ale změnit význam programu.}
        \end{itemize}

    \end{frame}

    \subsection{Lokální a globální optimalizace}

    \begin{frame}{Lokální a~globální optimalizace}

        \begin{itemize}
            \visible<2->{\item Některé optimalizace pracují pouze s krátkým úsekem kódu.}
            \visible<3->{\item Takovému úseku se často říká \textbf{základní blok}.}

            \visible<4->{\item \textbf{Lokální optimalizace} řeší jeden základní blok.}
            \visible<5->{
                \begin{itemize}
                    \item constant folding,
                    \item propagace kopií,
                    \item odstranění mrtvého kódu v bloku,
                    \item odstranění společných podvýrazů.
                \end{itemize}
            }

            \visible<6->{\item \textbf{Globální optimalizace} potřebují znát tok řízení programu.}
            \visible<7->{
                \begin{itemize}
                    \item analýza živých proměnných,
                    \item optimalizace přes větve programu,
                    \item optimalizace cyklů.
                \end{itemize}
            }
        \end{itemize}

    \end{frame}

    \subsection{Constant folding}

    \begin{frame}{Constant folding}

        \begin{itemize}
            \visible<2->{\item Pokud výraz závisí pouze na konstantách,
            může jej překladač spočítat už při překladu.}

            \visible<3->{\item Program za běhu potom nemusí provádět zbytečný výpočet.}
        \end{itemize}

        \vspace{0.3cm}

        \visible<4->{
            \centering
            \begin{tikzpicture}[
                box/.style={
                    draw,
                    rounded corners,
                    thick,
                align=center,
                font=\scriptsize,
                minimum width=2.0cm,
                minimum height=0.6cm
            }
        ]
            \node[box] (src) at (0,0) {\texttt{x := 2 * 60}\\ \texttt{y := x + 5}};
            \node[box] (res) at (6,0) {\texttt{y := 125}};
            \draw[->, thick] (src) -- (res);
        \end{tikzpicture}}

    \end{frame}

    \subsection{Mrtvý kód a společné podvýrazy}

    \begin{frame}{Mrtvý kód a společné podvýrazy}

        \visible<2->{\textbf{Odstranění mrtvého kódu}

        \[
        \begin{aligned}
            \mathtt{t1}\ &\mathtt{:= a * b}\\
            \mathtt{t2}\ &\mathtt{:= c + d}\\
            \mathtt{x}\  &\mathtt{:= t1}
        \end{aligned}
        \qquad
        \Longrightarrow
        \qquad
        \begin{aligned}
            \mathtt{t1} &\mathtt{:= a * b}\\
            \mathtt{x}  &\mathtt{:= t1}
        \end{aligned}
        \]}

        \vspace{0.25cm}

        \visible<3->{
            \textbf{Společný podvýraz}

            \[
            \begin{aligned}
                \mathtt{t1} &\mathtt{:= a * b}\\
                \mathtt{x}  &\mathtt{:= t1 + c}\\
                \mathtt{t2} &\mathtt{:= a * b}\\
                \mathtt{y}  &\mathtt{:= t2 + d}
            \end{aligned}
            \quad
            \Longrightarrow
            \quad
            \begin{aligned}
                \mathtt{t1} &\mathtt{:= a * b}\\
                \mathtt{x}  &\mathtt{:= t1 + c}\\
                \mathtt{y}  &\mathtt{:= t1 + d}
            \end{aligned}
            \]
        }

    \end{frame}

    \begin{frame}{Globální optimalizace}

        \begin{itemize}
            \visible<2->{\item Globální optimalizace už nestačí provádět v jednom bloku.}
            \visible<3->{\item Překladač musí znát \textbf{tok řízení programu}.}
            \visible<4->{\item Program se proto často reprezentuje jako graf základních bloků.}
        \end{itemize}

        \vspace{0.1cm}

        \centering
        \begin{tikzpicture}

            \node<4->[task] (b1) at (0,0) {$B_1$};
            \node<4->[task] (b2) at (-1.8,-1.2) {$B_2$};
            \node<4->[task] (b3) at (1.8,-1.2) {$B_3$};
            \node<4->[task] (b4) at (0,-2.4) {$B_4$};

            \draw<4->[diredge] (b1) -- (b2);
            \draw<4->[diredge] (b1) -- (b3);
            \draw<4->[diredge] (b2) -- (b4);
            \draw<4->[diredge] (b3) -- (b4);

        \end{tikzpicture}

        \vspace{0.1cm}

        \visible<5->{
            \begin{center}
                \scriptsize
                Globální optimalizace řeší například živost proměnných,
                optimalizace přes větve programu a optimalizace cyklů.
            \end{center}
        }

    \end{frame}

    \begin{frame}{Živé a mrtvé bloky programu}

        \begin{itemize}
            \visible<2->{\item Hrana \((B_i, B_j)\) znamená, že po vykonání bloku \(B_i\)
            může program pokračovat blokem \(B_j\).}
            \visible<3->{\item Blok je \textbf{živý}, pokud je dosažitelný z počátečního bloku.}
            \visible<4->{\item Blok, do kterého se program nikdy nemůže dostat, je \textbf{mrtvý}.}
        \end{itemize}

        \vspace{0.05cm}

        \visible<5->{\scriptsize
        \begin{center}
            \begin{minipage}{0.9\textwidth}
                \begin{algorithm}[H]
                    \caption{Detekce mrtvých bloků}
                    \KwIn{Graf toku řízení \(G=(B,E)\), počáteční blok \(B_0\)}
                    \KwOut{Rozdělení bloků na živé a mrtvé}

                    \ForEach{\textup{blok \(B_i\in B\)}}{
                        $\operatorname{stav}(B_i)\gets\textit{mrtvý}$
                    }

                    Založ zásobník \(S\) a vlož do něj \(B_0\)\\
                    $\operatorname{stav}(B_0)\gets\textit{živý}$\\

                    \While{\textup{zásobník \(S\) není prázdný}}{
                        Odeber blok \(B_i\) ze zásobníku\\

                        \ForEach{\textup{hrana \(B_i\to B_j\)}}{
                            \If{$\operatorname{stav}(B_j)=\textit{mrtvý}$}{
                                $\operatorname{stav}(B_j)\gets\textit{živý}$\\
                                Vlož \(B_j\) do zásobníku
                            }
                        }
                    }
                \end{algorithm}
            \end{minipage}
        \end{center}}

    \end{frame}

    \begin{frame}{Analýza živých proměnných}

        \begin{itemize}
            \visible<2->{\item Proměnná je \textbf{živá}, pokud se její hodnota může ještě použít.}
            \visible<3->{\item Pokud hodnota proměnné už nikdy nebude potřeba,
            lze některá přiřazení odstranit.}
        \end{itemize}

        \vspace{0.1cm}

        \visible<4->{\scriptsize
        \begin{center}
            \begin{minipage}{0.9\textwidth}
                \begin{algorithm}[H]
                    \caption{Zjednodušená analýza živých proměnných}
                    \KwIn{Základní blok instrukcí}
                    \KwOut{Množiny živých proměnných před a za instrukcemi}

                    $L\gets$ proměnné živé na konci bloku\\

                    \ForEach{\textup{instrukce $I$ procházené odzadu}}{
                        \tcp{\(\operatorname{liveAfter}(I)\) jsou proměnné, jejichž hodnoty mohou být potřeba po vykonání instrukce \(I\).}
                        $\operatorname{liveAfter}(I)\gets L$\\

                        Odstraň z $L$ proměnnou, do které instrukce $I$ zapisuje\\
                        Přidej do $L$ proměnné, které instrukce $I$ čte\\

                        \tcp{\(\operatorname{liveBefore}(I)\) jsou proměnné, jejichž hodnoty musí být známy před vykonáním instrukce \(I\).}
                        $\operatorname{liveBefore}(I)\gets L$
                    }
                \end{algorithm}
            \end{minipage}
        \end{center}}

    \end{frame}

    \section{Interpretované jazyky}

    \titleframe{Interpretované jazyky}

    \subsection{Kompilovaný vs. interpretovaný překlad}

    \begin{frame}{Kompilovaný a interpretovaný překlad}

        \begin{columns}[T,onlytextwidth]
            \begin{column}{0.49\textwidth}
                \centering
                \textbf{Kompilovaný překlad}

                \vspace{0.1cm}

                \begin{tikzpicture}[
                    scale=0.75,
                    transform shape,
                    box/.style={
                        draw,
                        thick,
                        minimum width=18mm,
                        minimum height=10mm,
                        align=center,
                        font=\scriptsize
                    },
                    arrow/.style={->, thick, >=stealth}
                ]
                    \node[box] (src) at (0,0) {zdrojový\\program};
                    \node[box] (comp) at (2.1,0) {překladač};
                    \node[box] (target) at (4.2,0) {cílový\\program};
                    \node[box] (data) at (4.2,1.25) {data};
                    \node[box] (out) at (6.3,0) {výsledky};

                    \draw[arrow] (src) -- (comp);
                    \draw[arrow] (comp) -- (target);
                    \draw[arrow] (target) -- (out);
                    \draw[arrow] (data) -- (target);
                \end{tikzpicture}
            \end{column}

            \begin{column}{0.49\textwidth}
                \centering
                \textbf{Interpretovaný překlad}

                \vspace{0.1cm}

                \begin{tikzpicture}[
                    scale=0.75,
                    transform shape,
                    box/.style={
                        draw,
                        thick,
                        minimum width=18mm,
                        minimum height=10mm,
                        align=center,
                        font=\scriptsize
                    },
                    arrow/.style={->, thick, >=stealth}
                ]
                    \node[box] (src) at (0,0) {zdrojový\\program};
                    \node[box] (interp) at (2.1,0) {interpret};
                    \node[box] (data) at (2.1,1.25) {data};
                    \node[box] (out) at (4.2,0) {výsledky};

                    \draw[arrow] (src) -- (interp);
                    \draw[arrow] (data) -- (interp);
                    \draw[arrow] (interp) -- (out);
                \end{tikzpicture}
            \end{column}
        \end{columns}

        \vspace{0.3cm}

        \begin{itemize}
            \visible<2->{\item U kompilovaného jazyka nejprve vznikne \textbf{cílový program}.}
            \visible<3->{\item U interpretovaného jazyka program vykonává přímo \textbf{interpret}.}
        \end{itemize}

    \end{frame}

    \subsection{Jak funguje interpret?}

    \begin{frame}{Jak funguje interpret?}

        \begin{itemize}
            \visible<2->{\item \textbf{Interpret} je program, který čte zdrojový program a rovnou jej vykonává.}
            \visible<3->{\item Obvykle postupuje po jednotlivých příkazech nebo instrukcích.}
            \visible<4->{\item Nepřekládá celý program předem do samostatného spustitelného souboru.}
        \end{itemize}

        \vspace{0.3cm}

        \centering
        \begin{tikzpicture}[
            box/.style={
                draw,
                rounded corners,
                thick,
                minimum width=22mm,
                minimum height=8mm,
                align=center,
                font=\scriptsize
            },
            arrow/.style={->, thick, >=stealth}
        ]
            \node<2->[box] (src) at (0,0) {zdrojový\\program};
            \node<3->[box] (read) at (2.5,0) {načtení\\příkazu};
            \node<4->[box] (exec) at (5.0,0) {vykonání};
            \node<5->[box] (next) at (7.5,0) {další\\příkaz};

            \draw<3->[arrow] (src) -- (read);
            \draw<4->[arrow] (read) -- (exec);
            \draw<5->[arrow] (exec) -- (next);
            \draw<5->[arrow] (next.south) to[out=-90,in=-90] (read.south);
        \end{tikzpicture}

    \end{frame}

    \subsection{Výhody a nevýhody}

    \begin{frame}{Výhody a nevýhody}

        \begin{columns}[T,onlytextwidth]
            \begin{column}{0.48\textwidth}
                \begin{timportantbox}{Výhody}
                    \begin{itemize}
                        \visible<2->{\item jednodušší spuštění programu,}
                        \visible<3->{\item dobrá přenositelnost mezi platformami,}
                        \visible<4->{\item snadnější ladění a testování,}
                        \visible<5->{\item rychlejší vývoj a úpravy programu.}
                    \end{itemize}
                \end{timportantbox}
            \end{column}

            \begin{column}{0.48\textwidth}
                \begin{talertbox}{Nevýhody}
                    \begin{itemize}
                        \visible<6->{\item obvykle pomalejší běh programu,}
                        \visible<7->{\item interpret musí být přítomen při spuštění,}
                        \visible<8->{\item část práce se opakuje při každém běhu.}
                    \end{itemize}
                \end{talertbox}
            \end{column}
        \end{columns}

    \end{frame}

    \begin{frame}{Co zpomaluje interpretované jazyky?}

        \begin{itemize}
            \visible<2->{\item Interpret obvykle vykonává program \textbf{postupně za běhu}.}

            \visible<3->{\item Při každém spuštění se může znovu opakovat část práce:}
            \visible<4->{
                \begin{itemize}
                    \item načítání a rozbor programu,
                    \item převod příkazů na vnitřní reprezentaci,
                    \item kontrola typů za běhu,
                    \item vyhledávání proměnných a funkcí,
                    \item rozhodování, jakou operaci právě provést.
                \end{itemize}
            }

            \visible<5->{\item Místo přímého vykonávání strojových instrukcí běží
            \textbf{interpretační smyčka}.}
        \end{itemize}

        \vspace{0.1cm}

        \onslide<6->
        \begin{tinfoframe}
            \textbf{Zjednodušeně:} U~kompilovaného programu mnoho rozhodnutí proběhne už při překladu.
            U~interpretovaného programu se část těchto rozhodnutí přesouvá až do běhu programu.
        \end{tinfoframe}

    \end{frame}

    \subsection{JIT kompilace}

    \begin{frame}{JIT kompilace}

        \begin{itemize}
            \visible<2->{\item \textbf{JIT} = \emph{Just-In-Time compilation}.}
            \visible<3->{\item Program se nepřekládá celý dopředu, ale \textbf{za běhu}.}
            \visible<4->{\item Často se nejprve vytvoří \textbf{mezikód} nebo \textbf{bytecode}.}
            \visible<5->{\item Ten je potom při spuštění překládán do strojového kódu.}
        \end{itemize}

        \vspace{0.3cm}

        \centering
        \begin{tikzpicture}[
            box/.style={
                draw,
                rounded corners,
                thick,
                minimum width=20mm,
                minimum height=8mm,
                align=center,
                font=\scriptsize
            },
            arrow/.style={->, thick, >=stealth}
        ]
            \node<3->[box] (src) at (0,0) {zdrojový\\kód};
            \node<4->[box] (bc) at (2.6,0) {bytecode\\/ mezikód};
            \node<5->[box] (jit) at (5.2,0) {JIT\\kompilátor};
            \node<6->[box] (native) at (7.8,0) {strojový\\kód};

            \draw<4->[arrow] (src) -- (bc);
            \draw<5->[arrow] (bc) -- (jit);
            \draw<6->[arrow] (jit) -- (native);
        \end{tikzpicture}

        \vspace{0.3cm}

        \onslide<7->
        \begin{timportantframe}
            \begin{center}
                JIT spojuje výhody interpretace a kompilace:
                \textbf{přenositelnost mezikódu} a \textbf{vyšší rychlost běhu}.
            \end{center}
        \end{timportantframe}

    \end{frame}

    \begin{frame}{Proč je dnes JIT preferovaná?}

        \begin{itemize}
            \visible<2->{\item Čistá interpretace je jednoduchá, ale často příliš pomalá.}
            \visible<3->{\item Klasická kompilace je rychlá při běhu, ale hůře se přizpůsobuje běhu programu (hlavně pro účely \textbf{debugování}).}
            \visible<4->{\item JIT dokáže přeložit jen ty části programu,
            které se skutečně často používají.}
            \visible<5->{\item Díky tomu může optimalizovat podle skutečného chování programu za běhu.}
        \end{itemize}

    \end{frame}

    \begin{frame}{Příklady jazyků a prostředí s JIT}

        \begin{itemize}
            \visible<2->{\item \textbf{Java} -- program se překládá do bytecode pro JVM,
            která používá JIT kompilaci.}

            \visible<3->{\item \textbf{C\#} -- kód se překládá do mezikódu IL,
            který běhové prostředí .NET překládá za běhu.}

            \visible<4->{\item \textbf{JavaScript} -- moderní prohlížeče
            (např. V8 nebo SpiderMonkey) používají JIT.}

            \visible<5->{\item \textbf{Python} -- klasická implementace CPython je spíše interpret,
            ale například \textbf{PyPy} používá JIT.}

            \visible<6->{\item \textbf{LuaJIT} -- speciální velmi rychlá implementace jazyka Lua.}
        \end{itemize}

        \vspace{0.1cm}

        \onslide<7->
        \begin{tinfoframe}
            Dnes se proto v praxi často nesetkáváme s "čistě interpretovanými"
            jazyky, ale spíše s běhovými prostředími, která využívají JIT kompilaci.
        \end{tinfoframe}

    \end{frame}

    \section{Zdroje}

    \begin{frame}{Zdroje}
        \begin{itemize}
            \item ČEŠKA, Milan, Tomáš HRUŠKA a Miroslav BENEŠ. \textit{Překladače.} Brno: Vysoké učení technické v Brně, Fakulta elektrotechnická, b. r. Učební texty vysokých škol. Dostupné \href{https://www.fi.muni.cz/usr/kretinsky/prekladace_skripta_VUT.pdf}{\underline{online}}. Cit. 2026-07-02.
            \item BRYANT, Randal E. a David R. O'HALLARON. \textit{Computer Systems: A Programmer's Perspective.} 3rd ed. Boston: Pearson, 2016. ISBN 978-0-13-409266-9.
        \end{itemize}
    \end{frame}

\end{document}